\begin{abstract}
Medidas de conservação de água podem, involuntariamente, elevar a salinidade do esgoto tratado ao reduzir o volume de água que dilui a mesma carga de sais. Este trabalho investiga a tendência de longo prazo de Sólidos Dissolvidos Totais (TDS) no efluente da Los Angeles--Glendale Water Reclamation Plant (LAGWRP), a partir de 182 meses de dados reais do portal eSMR (2011--2026). Uma bateria de 21 métodos de previsão -- estatística clássica de tendência, séries temporais (STL/SARIMA), modelos baseados em árvore, Prophet, regressão bayesiana, SVR, Gaussian Process, tendência+árvore no resíduo, SARIMAX multivariado com Cloreto, ensemble SARIMA+Prophet, LSTM leve e quatro \textit{baselines} obrigatórios -- foi avaliada sob um mesmo protocolo de validação temporal honesta (MASE, validação cruzada expansiva, \textit{backtesting} de origem móvel), com diagnóstico de resíduos dos candidatos mais fortes. Os métodos com ajuste estatístico global convergem para uma tendência de alta estatisticamente significativa de 2,9 a 3,9 mg/L/ano; três finalistas complementares (regressão bayesiana, tendência+Random Forest no resíduo, e um ensemble SARIMA+Prophet) projetam entre ~760 e ~800 mg/L em +20 anos, apesar de mecanismos de extrapolação distintos. TDS correlaciona-se positivamente com Amônia mesmo após remoção de tendência comum ($r=0{,}175$, p=0,018), consistente com a hipótese de inibição da nitrificação por salinidade; a correlação com BOD não é estatisticamente significativa, o que parece refletir a baixa variação real dessa série (rente ao limite de detecção do método) mais do que ausência de efeito. Os achados são compatíveis com o mecanismo de conservação de água descrito na literatura consultada.
\end{abstract}

\noindent \textbf{Palavras-chave:} Salinidade de efluente, séries temporais, previsão, validação temporal, tratamento de esgoto.
